\Article{5}
The first proof of the theorem is due to the distinguished mathematician \textsc{d'Alembert}, \textit{Recherches sur le calcul intégral, Histoire de l'Acad. de Berlin, Année} 1746, p. 182 ff. It is also found in \textsc{Bougainville}, \textit{Traité du calcul intégral, à Paris} 1754, p. 47 ff. The principal steps of this method are as follows.

First he shows that, if any function $X$ of the variable $x$ becomes $=0$ either for $x=0$ or for $x=\infty$, and can attain an infinitely small positive real value by assigning a real value to $x$, this function can also attain an infinitely small negative real value through a value of $x$ that is either real or of the imaginary form $p+q\sqrt{-1}$. Namely, if $\Omega$ denotes an infinitely small value of $X$, and $\omega$ the corresponding value of $x$, he asserts that $\omega$ can be expressed by a rapidly convergent series $a\Omega^\alpha+b\Omega^\beta+c\Omega^\gamma$, etc., where the exponents $\alpha$, $\beta$, $\gamma$, etc. are continually increasing rational quantities, which therefore, at least beyond some fixed distance from the beginning, become positive and make the terms in which they occur infinitely small. Now, if none of these exponents is a fraction with an even denominator, all the terms of the series become real for both positive and negative values of $\Omega$. If, however, there are fractions with even denominators among these exponents, then for a negative value of $\Omega$ the corresponding terms are known to have the form $p+q\sqrt{-1}$. But because of the infinite convergence of the series, in the former case it is enough to retain only the first (that is, the largest) term; in the latter there is no need to proceed beyond the term that first produces an imaginary part.

By similar reasoning it can be shown that, if $X$ can attain an infinitely small negative value from a real value of $x$, this function can attain an infinitely small positive real value from a real value of $x$ or an imaginary one of the form $p+q\sqrt{-1}$.

From this he concludes, secondly, that there is also some finite real value of $X$, negative in the former case and positive in the latter, that can be produced by an imaginary value of $x$ of the form $p+q\sqrt{-1}$.

\SourceLeaf{013}
It follows that, if $X$ is a function of $x$ that attains a real value $V$ from a real value $v$ of $x$, and also attains from a real value of $x$ a real value greater or smaller by an infinitely small quantity, it can also take a real value smaller or greater, respectively, than $V$ by an infinitely small and hence by a finite quantity, by assigning to $x$ a value of the form $p+q\sqrt{-1}$. This is readily derived from what precedes by conceiving $V+Y$ substituted for $X$, and $v+y$ for $x$.

Finally, the distinguished \textsc{d'Alembert} asserts that, if $X$ is supposed capable of traversing an entire interval between two real values $R,S$ (that is, of becoming equal both to $R$ and to $S$, and to all intermediate real values), by assigning to $x$ values always of the form $p+q\sqrt{-1}$, the function $X$ can be increased or decreased still further by any finite real quantity (according as $S>R$ or $S<R$), while $x$ always remains of the form $p+q\sqrt{-1}$. For if there were a real quantity $U$ (with $S$ supposed to lie between it and $R$) to which $X$ could not become equal through such a value of $x$, there must be a \textit{maximum} value of $X$ (when $S>R$; a minimum when $S<R$), say $T$, attained from a value $p+q\sqrt{-1}$ of $x$, such that no value of the same form could be assigned to $x$ that would bring the function $X$ any closer to $U$, even by the smallest amount. Now, if in this equation between $X$ and $x$ we substitute $p+q\sqrt{-1}$ everywhere for $x$ and then set equal to zero both the real part and the part containing the factor $\sqrt{-1}$, after omitting that factor, two other equations can be obtained by elimination from the two resulting equations (which contain $p,q$ and $X$ together with constants). One contains $p,X$ and constants; the other, free of $p$, contains only $q,X$ and constants. Thus, since $X$ has traversed all values from $R$ to $T$ through real values of $p,q$, what precedes shows that $X$ can approach $U$ still more closely by assigning to $p,q$ values of the forms $\alpha+\gamma\sqrt{-1}$, $\beta+\delta\sqrt{-1}$, respectively. But then $x=\alpha-\delta+(\gamma+\beta)\sqrt{-1}$, that is, it is still of the form $p+q\sqrt{-1}$, contrary to the hypothesis.

Now, if $X$ is supposed to denote a function such as $x^m+Ax^{m-1}+Bx^{m-2}+\text{etc.}+M$, it is readily seen that real values can be assigned to $x$ such that $X$ traverses an entire interval between two real values. Hence $x$ can attain some value of the form $p+q\sqrt{-1}$ from which $X$ becomes $=0$. Q. E. D.\HistoricalNote{2}
\SourceLeaf{014}
\Article{6}
The objections that seem possible against \textsc{d'Alembert}'s proof come down chiefly to the following.

1. The distinguished d'A. raises no doubt about the \textit{existence} of values of $x$ corresponding to given values of $X$, but assumes it and investigates only the \textit{form} of those values.

Although this objection is in itself very serious, here it concerns only the form of expression, which can easily be corrected so as to remove it entirely.

2. The assertion that $\omega$ can always be expressed by a series of the kind he gives is certainly false if $X$ is also to designate any transcendental function (as d'A. suggests in several places). This is evident, for example, if we put $X=e^{\frac1x}$, or $x=\frac1{\log X}$. Yet if we restrict the proof to the case in which $X$ is an algebraic function of $x$ (which suffices for the present purpose), the proposition is indeed true. --- d'A., moreover, has offered nothing to support his supposition. The celebrated \textsc{Bougainville} supposes $X$ to be an algebraic function of $x$ and recommends \textsc{Newton}'s parallelogram for finding the series.

3. He uses infinitely small quantities more freely than is consistent with geometric rigour, or than a careful analyst would allow, at least in our own time (when they have a deservedly bad reputation). Nor does he make the leap from an infinitely small value of $\Omega$ to a finite one sufficiently clear. He seems to infer his proposition that $\Omega$ can also attain some finite value, less from the possibility of an infinitely small value of $\Omega$ than from this: if $\Omega$ denotes a very small quantity, then because of the rapid convergence of the series, the more terms of the series are taken, the closer one comes to the true value of $\omega$; or, the more terms are included in the sum taken for $\omega$, the more accurately the equation expressing the relation between $\omega$ and $\Omega$, or $x$ and $X$, will be satisfied. Apart from the fact that this entire argument seems too vague to yield any rigorous conclusion, I note that there are indeed series that, however small a value is assigned to the quantity
\SourceLeaf{015}
in whose powers they proceed, nevertheless always diverge, so that, if only they are continued far enough, one can reach terms greater than any given quantity\HistoricalNote{3}. This happens when the coefficients of the series form a hypergeometric progression. It was therefore necessary to prove that such a hypergeometric series could not arise in the present case.

Moreover, it seems to me that the distinguished d'A. was wrong to resort to infinite series here, and that they are unsuitable for establishing this fundamental theorem of the theory of equations.

4. From the supposition that $X$ can attain the value $S$, but not the value $U$, it does not yet follow that between $S$ and $U$ there must lie a value $T$ that $X$ can reach but cannot exceed. Another case still remains: there could be a limit between $S$ and $U$ that $X$ can approach as closely as one wishes, yet never reach. From the arguments given by the distinguished d'A. it follows only that $X$ can exceed every value it has reached by a further finite quantity. For instance, once it has become $=S$, it can be increased by some further finite quantity $\Omega$; after this a new increment $\Omega'$ can be added, then another increment $\Omega''$, etc., so that however many increments have already been added, none should be regarded as the last, but a new one can always be added. Yet although the \textit{number} of possible increments is bounded by no limits, it could certainly happen that, if the increments $\Omega,\Omega',\Omega''$, etc. continually decreased, the sum $S+\Omega+\Omega'+\Omega''$, etc. would nevertheless never attain some limit, however many terms were considered.

Although this case cannot occur when $X$ denotes a polynomial in $x$, without a proof that it cannot happen the method must nevertheless be regarded as incomplete. When, however, $X$ is a transcendental function, or even an algebraic fractional function, that case can indeed occur: for example, whenever an infinitely large value of $x$ corresponds to some value of $X$\SourceLeaf{016}. Then \textsc{d'Alembert}'s method seems incapable of being reduced to unquestionable principles without many complications, and perhaps in some cases not at all.

For these reasons I cannot regard \textsc{d'Alembert}'s proof as satisfactory. Nevertheless, it seems to me that its real force is by no means broken by all these objections. I believe that on the same foundation (though in a very different way, and at least with greater circumspection) not only can a rigorous proof of our theorem be built, but everything one could wish for in the theory of \textit{transcendental} equations can be obtained from it. I shall discuss this very important matter more fully on another occasion; meanwhile, compare Article 24 below.
